\documentclass[11pt]{article}
% ===========================================================================
%  Shared preamble for the AP Calculus AB exam-preparation worksheets.
%  Input right after \documentclass. Files beginning with "_" are not compiled
%  on their own. Navy and gold, to match the course's lesson pages.
% ===========================================================================
\usepackage[letterpaper,top=2.3cm,bottom=2.2cm,left=1.7cm,right=1.7cm,headheight=28pt,footskip=30pt]{geometry}
\usepackage{amsmath,amssymb}
\usepackage[T1]{fontenc}
\usepackage[utf8]{inputenc}
\usepackage{lmodern}
\usepackage{xcolor}
\usepackage[most]{tcolorbox}
\usepackage{enumitem}
\usepackage{fancyhdr}
\usepackage{array}
\usepackage{tikz}
\usepackage{pgfplots}
\pgfplotsset{compat=1.18}

\definecolor{navy}{HTML}{161E54}
\definecolor{ink}{HTML}{27357E}
\definecolor{gold}{HTML}{B98A00}
\definecolor{goldbg}{HTML}{FFFDF3}
\definecolor{inkbg}{HTML}{F6F7FD}
\definecolor{rule}{HTML}{C9CFF0}
\definecolor{okgreen}{HTML}{1B7A44}
\definecolor{okbg}{HTML}{F1F8F3}

\pagestyle{fancy}
\fancyhf{}
\fancyhead[L]{\footnotesize Name: \underline{\hspace{4cm}}}
\fancyhead[C]{\textbf{\color{navy}Integration Academy}\\[-2pt]{\scriptsize Calculus AB \textperiodcentered{} AP\textsuperscript{\textregistered} Exam Preparation}}
\fancyfoot[L]{\tiny\color{black!55}\parbox[t]{0.78\textwidth}{AP\textsuperscript{\textregistered} is a registered trademark of the College Board, which was not involved in the production of, and does not endorse, this material. \textcopyright{} Integration Academy.}}
\fancyfoot[R]{\scriptsize\thepage}
\renewcommand{\headrulewidth}{1.2pt}
\renewcommand{\headrule}{\hbox to\headwidth{\color{gold}\leaders\hrule height \headrulewidth\hfill}}
\renewcommand{\footrulewidth}{0pt}
\setlength{\parindent}{0pt}
\setlength{\parskip}{4pt}

% Title block: \wstitle{code}{title}{unit line}
\newcommand{\wstitle}[3]{\begin{center}{\small\bfseries\color{gold}WORKSHEET #1}\\[3pt]{\LARGE\bfseries\color{navy}#2}\\[3pt]{\small\color{black!60}#3}\end{center}\vspace{2pt}}
% Section heading
\newcommand{\sech}[1]{\par\vspace{10pt}{\large\bfseries\color{navy}#1}\par\nobreak\vspace{2pt}{\color{rule}\hrule height0.8pt}\vspace{7pt}}
% A level tag printed beside a group of questions
\newcommand{\level}[1]{\par\vspace{6pt}{\footnotesize\bfseries\color{gold}\MakeUppercase{#1}}\par\nobreak\vspace{2pt}}

\newtcolorbox{ideabox}{colframe=ink,colback=inkbg,coltitle=white,colbacktitle=ink,fonttitle=\bfseries\small,title={Key ideas},arc=2pt,boxrule=0.8pt,left=8pt,right=8pt}
\newtcolorbox{example}[1]{colframe=rule,colback=white,coltitle=navy,colbacktitle=inkbg,fonttitle=\bfseries\small,title={#1},arc=2pt,boxrule=0.8pt,borderline west={3pt}{0pt}{ink},left=10pt,right=8pt}
\newtcolorbox{frbox}[1]{colframe=gold,colback=goldbg,coltitle=white,colbacktitle=gold,fonttitle=\bfseries\small,title={#1},arc=2pt,boxrule=0.8pt,breakable,left=8pt,right=8pt}
\newtcolorbox{solbox}{colframe=okgreen,colback=okbg,coltitle=white,colbacktitle=okgreen,fonttitle=\bfseries\small,title={Worked solutions},arc=2pt,boxrule=0.8pt,breakable,left=8pt,right=8pt}
\newtcolorbox{scorebox}[1]{colframe=navy,colback=white,coltitle=white,colbacktitle=navy,fonttitle=\bfseries\small,title={#1},arc=2pt,boxrule=0.8pt,breakable,left=8pt,right=8pt}

\newcommand{\soln}{\par\smallskip\textbf{\color{ink}Solution.}\ }
\newcommand{\result}[1]{\par\smallskip\textbf{#1}}
% blank working space of a given height
\newcommand{\work}[1]{\par\vspace{#1}}
% a point value in the margin of a scoring guide
\newcommand{\pts}[1]{\hfill{\footnotesize\bfseries\color{gold}[#1]}}

% Coordinate grid for graph questions: \begin{apgraph}{xmin}{xmax}{ymin}{ymax} ... \end{apgraph}
\newenvironment{apgraph}[4]{\begin{center}\begin{tikzpicture}\begin{axis}[
  axis lines=middle, xmin=#1, xmax=#2, ymin=#3, ymax=#4,
  grid=both, grid style={black!12}, xtick distance=1, ytick distance=1,
  ticklabel style={font=\scriptsize}, xlabel={$x$}, ylabel={$y$},
  width=10cm, height=7.2cm, samples=80, clip=false,
  every axis plot/.append style={line width=1.1pt, color=ink}]}{\end{axis}\end{tikzpicture}\end{center}}
\newcommand{\closeddot}[2]{\addplot[only marks, mark=*, mark size=2.2pt, color=ink] coordinates {(#1,#2)};}
\newcommand{\opendot}[2]{\addplot[only marks, mark=*, mark size=2.2pt, color=ink, mark options={fill=white}] coordinates {(#1,#2)};}

% Multiple-choice options: four across, or stacked when the options are long.
\newcommand{\opts}[4]{\par\vspace{3pt}\begin{tabular}{@{}p{0.235\linewidth}p{0.235\linewidth}p{0.235\linewidth}p{0.235\linewidth}@{}}(A)\ #1 & (B)\ #2 & (C)\ #3 & (D)\ #4\end{tabular}\par\vspace{9pt}}
\newcommand{\optsv}[4]{\par\vspace{3pt}\begin{tabular}{@{}l@{\ \ }p{0.9\linewidth}@{}}(A) & #1\\(B) & #2\\(C) & #3\\(D) & #4\end{tabular}\par\vspace{9pt}}

\fancyhead[R]{\footnotesize\color{black!60}Worksheet 2.1}
\begin{document}

\wstitle{2.1}{Slope at a Single Point}{Unit 2: The Derivative}

\begin{ideabox}
\begin{itemize}[leftmargin=*,itemsep=2pt,topsep=2pt]
  \item The \textbf{average rate of change} of $f$ on $[a,a+h]$ is the slope of a secant line: $\dfrac{f(a+h)-f(a)}{h}$.
  \item The \textbf{instantaneous rate of change} at $x=a$ is the limit of those slopes as $h\to0$. It is the slope of the tangent line.
  \item The tangent line at $x=a$ passes through $\big(a,f(a)\big)$: \ $y-f(a)=m(x-a)$.
  \item A rate of change has units of \textquotedblleft{}output per unit of input\textquotedblright{}.
\end{itemize}
\end{ideabox}

\sech{Worked examples}
\begin{example}{Example 1 \textperiodcentered{} Slope and tangent line}
Find the slope of the tangent line to $f(x)=x^2+1$ at $x=3$, and write its equation.\soln $\dfrac{(3+h)^2+1-10}{h}=\dfrac{6h+h^2}{h}=6+h\to6$. The point is $(3,10)$.\result{Slope 6; tangent line $y-10=6(x-3)$.}
\end{example}

\begin{example}{Example 2 \textperiodcentered{} A reciprocal}
Find the instantaneous rate of change of $f(x)=\dfrac4x$ at $x=2$.\soln $\dfrac{\frac{4}{2+h}-2}{h}=\dfrac{4-2(2+h)}{h(2+h)}=\dfrac{-2}{2+h}\to-1$.\result{The rate of change is $-1$.}
\end{example}

\begin{example}{Example 3 \textperiodcentered{} Velocity}
A ball has height $H(t)=50t-5t^2$ metres after $t$ seconds. Find its velocity at $t=2$.\soln $\dfrac{H(2+k)-H(2)}{k}=\dfrac{50k-20k-5k^2}{k}=30-5k\to30$.\result{At $t=2$ the ball is rising at 30 metres per second.}
\end{example}

\newpage
\sech{Practice}

\level{Core}
\begin{enumerate}[leftmargin=*,itemsep=2pt]
  \item Find the average rate of change of $f(x)=x^2-4$ on $[1,3]$.\work{2.3cm}
  \item Find the slope of the tangent line to $f(x)=3x^2$ at $x=1$.\work{2.3cm}
  \item Write an equation of the tangent line to $f(x)=x^2-2x$ at $x=4$.\work{2.3cm}
  \item Explain why the tangent line to $f(x)=5x-7$ has slope 5 at every point.\work{2.3cm}
\end{enumerate}

\level{Exam level}
\begin{enumerate}[leftmargin=*,itemsep=2pt,start=5]
  \item Find the slope of the tangent line to $f(x)=\dfrac{2}{x-1}$ at $x=3$.\work{2.8cm}
  \item Find the slope of the tangent line to $f(x)=\sqrt{x+5}$ at $x=4$.\work{2.8cm}
  \item Find the slope of the tangent line to $f(x)=x^3$ at $x=2$, and write its equation.\work{2.8cm}
  \item At what value of $x$ does the graph of $f(x)=x^2-6x$ have a horizontal tangent line? Use a limit to find the slope at a general point $x$.\work{2.8cm}
\end{enumerate}

\begin{frbox}{Question 9 \textperiodcentered{} Free-response practice}
A cart moves along a track so that its position after $t$ seconds is $s(t)=t^2+4t$ metres.
\begin{enumerate}[label=(\alph*),leftmargin=*,itemsep=2pt]
  \item Find the average velocity of the cart on the interval $[1,3]$.\work{2.6cm}
  \item For $h\neq0$, simplify the average velocity on $[2,2+h]$, and use it to find the velocity at $t=2$.\work{2.6cm}
  \item Compare your answers to parts (a) and (b). Is this a coincidence? Explain using your expression from part (b).\work{2.6cm}
\end{enumerate}
\end{frbox}

\newpage
\begin{solbox}
\begin{enumerate}[leftmargin=*,itemsep=4pt]
  \item $\dfrac{5-(-3)}{2}=4$.
  \item $\dfrac{3(1+h)^2-3}{h}=6+3h\to6$.
  \item $\dfrac{(4+h)^2-2(4+h)-8}{h}=6+h\to6$; the point is $(4,8)$: $y-8=6(x-4)$.
  \item $\dfrac{5(a+h)-7-(5a-7)}{h}=5$ for every $a$ and every $h$: a line is its own tangent.
  \item $\dfrac{\frac{2}{2+h}-1}{h}=\dfrac{-1}{2+h}\to-\dfrac12$.
  \item $\dfrac{\sqrt{9+h}-3}{h}=\dfrac{1}{\sqrt{9+h}+3}\to\dfrac16$.
  \item $\dfrac{(2+h)^3-8}{h}=12+6h+h^2\to12$; tangent line $y-8=12(x-2)$.
  \item $\dfrac{(x+h)^2-6(x+h)-x^2+6x}{h}=2x+h-6\to2x-6$, which is zero at $x=3$.
  \item \begin{enumerate}[label=(\alph*),leftmargin=*,itemsep=3pt]
    \item $\dfrac{s(3)-s(1)}{2}=\dfrac{21-5}{2}=8$ metres per second.
    \item $\dfrac{(2+h)^2+4(2+h)-12}{h}=\dfrac{8h+h^2}{h}=8+h$, which approaches 8 metres per second.
    \item They are equal, and not by coincidence. On $[2-k,2+k]$ the average velocity is $\dfrac{s(2+k)-s(2-k)}{2k}=\dfrac{16k}{2k}=8$: for a quadratic, the average velocity over an interval equals the velocity at its midpoint.
  \end{enumerate}
\end{enumerate}
\end{solbox}

\end{document}
